\section{Two Designs}\label{sec:designs}

\Cref{tab:designs} sets the two designs side by side. Both are \m{model} at a
different triple, and that difference is the whole security story.

\begin{table*}[t]
\centering
\caption{The two designs.}
\label{tab:designs}
\footnotesize
\begin{tabularx}{\textwidth}{@{}lLL@{}}
\toprule
 & \m{model_immediate} & \m{model_sliced} \\
\midrule
\code{init} & all masks clear; no slice counter & every mask set except the timer's;
  counter at \code{Some 0} \\
\code{handler_preroutine} & on a handler's \code{Notify}, unmask everything; on the timer's,
  also request a reschedule & on the timer's \code{Notify}, open a slice and request a
  reschedule; at each boundary inside it, reopen the masks and queue a NOP run; tick
  the slice and close it \\
\m{restore_invariant} & \code{id} & re-imposes the slice invariant; the identity on every
  reachable state \\
a handler runs & as soon as its interrupt is serviced & only inside a slice opened
  by the public timer \\
non-interfering? & no (\cref{thm:leak}) & yes (\cref{thm:main}) \\
\bottomrule
\end{tabularx}
\end{table*}

\subsection{Ordinary interrupt handling, and the leak}\label{sec:leak}

\m{model_immediate} starts with every mask clear, so an interrupt is serviced the
moment it arrives. Its preroutine reacts to a handler announcing completion by
unmasking everything, and if the handler was the timer's, by requesting a
reschedule. Nothing constrains \emph{when} a handler runs.

The counterexample is stated for a concrete system: a public process that
repeatedly emits the request \code{GetRequest}, a private process that issues a system
call when the disk handler has notified it and \code{NOP} otherwise, a round-robin
scheduler, and handlers of length two.

\begin{theorem}[\m{model_immediate_not_NI}]\label{thm:leak}
\begin{multline*}
\neg\,\NI\;\texttt{in\_equiv}\;\texttt{out\_equivC}\\
\texttt{model\_immediate\_concrete}.
\end{multline*}
\end{theorem}

\begin{proof}[Proof sketch]
Start from the trace of two public requests, which the model admits. Insert a disk
interrupt at the front; it is unobservable at $\bot$, so non-interference demands
that the result still be a trace. It is not. The interrupt sets the disk pending bit,
the first request still goes through, and the following step belongs to the disk
handler, so the public slot carries \code{None} where the second request was. Refuting
this requires showing that a trace is \emph{not} admitted, which needs reductions to
be invertible (\cref{sec:mechanisation}).
\end{proof}

\subsection{Sliced interrupt handling}\label{sec:sliced}

\m{model_sliced} keeps the disk and NOP handlers masked outside a slice. Its
preroutine is a composition of three stages, read right to left:
\begin{lstlisting}
sliced_preroutine o =
  check_ir_count (@$\circ$@) check_handler_completed
  (@$\circ$@) initiate_ir o
\end{lstlisting}
\begin{itemize}
\item \emph{Opening a slice.} When the timer handler emits \code{Notify}, \m{initiate_ir} loads
  the counter with $\mathit{runs}\cdot\mathit{runtime}$, so that exactly \code{runs}
  handler runs fit in the slice, and requests a reschedule for when the slice ends.
\item \emph{Boundaries.} Because the slice length is a multiple of \code{runtime}, the counter
  is a nonzero multiple of \code{runtime} exactly when a handler run has just finished.
  There \m{check_handler_completed} reopens the disk and NOP masks, keeps the timer
  masked, and makes the NOP handler pending. That last step is the \emph{filler}: a
  handler is always ready at a boundary, so a run takes place whether or not a disk
  interrupt is waiting, and the disk handler, if pending, takes priority.
\item \emph{Closing.} \m{check_ir_count} ticks the counter, and when it reaches zero masks the
  disk and NOP handlers, unmasks the timer and disables the counter.
\end{itemize}

A secret handler run therefore takes the place of a NOP run instead of displacing a
user process: a slice containing a disk run has the same length and position as one
containing only filler, and the public schedule around it is unchanged.
\Cref{fig:traces} shows two runs of the concrete system. In the run with a disk
interrupt, the interrupt is merely recorded until the next slice opens, and the only
structural difference between the runs is which handler occupies the first run of
the slice.

\begin{figure*}[t]
\centering
\footnotesize
\begin{tabular}{@{}lll@{}}
\toprule
no disk interrupt & with disk interrupt & \\
\midrule
\code{Get} & \code{Get} & \\
 & \textit{disk interrupt} & recorded as pending; nothing else happens \\
 & \code{Get}, \code{Get} & \\
\textit{timer interrupt} & \textit{timer interrupt} & \\
\code{Get} & \code{Get} & \\
\code{tmr Nth}, \code{tmr Nfy} & \code{tmr Nth}, \code{tmr Nfy} & timer handler completes: slice opens \\
\code{dfl Nth}, \code{dfl Nfy} & \code{dsk Nth}, \code{dsk Nfy} & the disk handler replaces a NOP run \\
\code{dfl Nth}, \code{dfl Nfy} & \code{dfl Nth}, \code{dfl Nfy} & \\
\code{sch hi} & \code{sch hi} & slice over: the scheduler runs \\
\code{NOP} & \code{Sys} & the private process reacts; both are secret \\
\bottomrule
\end{tabular}
\caption{Two runs of the concrete sliced system (abridged). Handler slots are
  secret including whether they ran, so to the attacker both columns look alike.}
\label{fig:traces}
\end{figure*}

\m{model_sliced} still exposes every slot of the pool. The model
\m{model_sliced_userview} applies a projection \m{parse_output} that keeps only the two
user processes' outputs; it is \m{model_sliced} behind a \code{map}, and its
non-interference follows by weakening the output equivalence.
